%!TEX TS-program = xelatex
%!TEX encoding = UTF-8 Unicode
% Mengting Li — Curriculum Vitae
% Built with Awesome-CV (https://github.com/posquit0/Awesome-CV), class file in this folder.
% Compile with:  xelatex resume.tex  (twice, for the "N of M" page footer)

\documentclass[11pt, a4paper]{awesome-cv}

\geometry{left=1.4cm, top=.8cm, right=1.4cm, bottom=1.8cm, footskip=.5cm}
\fontdir[fonts/]

\colorlet{awesome}{awesome-red}
\setbool{acvSectionColorHighlight}{true}
\renewcommand{\acvHeaderSocialSep}{\quad\textbar\quad}

\usepackage{enumitem}

%-------------------------------------------------------------------------------
%  PERSONAL INFORMATION
%-------------------------------------------------------------------------------
\name{Mengting}{Li}
\position{PhD Student in Bioinformatics}
\address{Room 3-08, The Hong Kong Jockey Club Building For Interdisciplinary Research, 5 Sassoon Rd, Sandy Bay, Hong Kong}
\pronouns{she/her}
\mobile{+852 5968 9107}
\email{mtlee1001@connect.hku.hk}
\homepage{limengting1001.com}
\github{huobei1997}
\linkedin{huobei}

\newcommand*{\monthyeartoday}{\ifcase\month\or January\or February\or March\or April\or May\or June\or July\or August\or September\or October\or November\or December\fi\ \number\year}

\Urlmuskip=0mu plus 1mu\relax
\def\UrlBreaks{\do\/\do-\do.\do:\do\_}

\begin{document}

\makecvheader
\makecvfooter{\monthyeartoday}{Mengting Li~~~·~~~Curriculum Vitae}{\thepage~of~\pageref{acvlastpage}}

%-------------------------------------------------------------------------------
%  EDUCATION
%-------------------------------------------------------------------------------
\cvsection{Education}

\begin{cventries}
  \cventry{PhD in Bioinformatics}{The University of Hong Kong (HKU)}{Hong Kong, China}{Oct 2025 - Present}{}
  \cventry{Master of Biological and Pharmaceutical Engineering}{University of Chinese Academy of Sciences (UCAS)}{Beijing, China}{Sep 2020 - Jun 2023}{}
  \cventry{Bachelor of Bioinformatics}{Xuzhou Medical University (XZMU)}{Xuzhou, China}{Sep 2016 - Jun 2020}{}
\end{cventries}

%-------------------------------------------------------------------------------
%  RESEARCH EXPERIENCE
%-------------------------------------------------------------------------------
\cvsection{Research Experience}

\begin{cvparagraph}
\textbf{Research Assistant (Jul 2023 - Sep 2025)}\\
\textbf{The University of Hong Kong | Hong Kong, China (Supervisor: Prof. Leo Poon)}\\
\textbf{Project: An interferon-stimulated long non-coding RNA USP30-AS1 as an immune modulator in influenza A virus infection}
\vspace{4.5mm}
\begin{cvitems}
  \item Developed a reproducible bioinformatics pipeline that successfully processed and analyzed 5 RNA-seq datasets and 4 microarray datasets from public repositories (SRA, GEO).
  \item Identified 1,715 differentially expressed lncRNAs across multiple post-infection time points by implementing a custom differential expression analysis workflow (using \texttt{STAR}, \texttt{featureCounts}, and \texttt{DESeq2}).
  \item Revealed key immune pathways using \texttt{g:Profiler} for functional enrichment analysis and \texttt{EnrichmentMap} in \texttt{Cytoscape} for network visualization.
\end{cvitems}

\vspace{4mm}
\textbf{Master's Student Researcher (Sep 2020 - Jun 2023)}\\
\textbf{BGI-Research \& University of Chinese Academy of Sciences | Shenzhen, China (Supervisor: Prof. Kui Wu)}\\
\textbf{Project: Multi-omics sequencing of gastroesophageal junction adenocarcinoma reveals prognosis-relevant key factors and a novel immunogenomic classification}
\vspace{4.5mm}
\begin{cvitems}
  \item Led a multi-omics analysis of 92 gastroesophageal junction adenocarcinoma (GEJAC) patients, integrating whole-genomic, transcriptomic, and T-cell receptor repertoires data.
  \item Developed a comprehensive bioinformatics pipeline for processing WGS data to perform somatic variant calling and driver gene identification. Further utilized these results for downstream analyses including mutational signature analysis (\texttt{SigProfilerExtraction}), neoantigen prediction (\texttt{NetMHC}, \texttt{NetMHCpan}), and tumor evolution (\texttt{PhylogicNDT}).
  \item Discovered 30 mutated driver genes, 7 novel genomic signatures, 3 copy-number variations, and 2 V-J gene usages previously unidentified and linked to prognosis.
  \item Developed a novel immunogenomic classification system that stratifies patients into three immune-related subtypes (immune-inflamed, intermediate, and deserted) with discrete profiles of genomic signatures, immune status, and clinical outcomes.
\end{cvitems}
\end{cvparagraph}

%-------------------------------------------------------------------------------
%  SKILLS
%-------------------------------------------------------------------------------
\cvsection{Skills}

\begin{cvparagraph}
\textbf{Programming \& Scripting Languages}
\vspace{4.5mm}
\begin{cvitems}
  \item \textbf{R}: (Proficient) Data manipulation (\texttt{Tidyverse}, \texttt{data.table}), statistical modeling, bioinformatics (\texttt{Bioconductor}), visualization (\texttt{ggplot2}).
  \item \textbf{Python}: (Proficient) Scientific computing (\texttt{Pandas}, \texttt{NumPy}), machine learning (\texttt{PyTorch}, \texttt{Scikit-learn}), scripting and workflow integration.
  \item \textbf{Bash/Shell Scripting}: Proficient in automating tasks, managing data, and working in a Linux/HPC environment.
\end{cvitems}

\vspace{4mm}
\textbf{Bioinformatics \& Genomics}
\vspace{4.5mm}
\begin{cvitems}
  \item \textbf{Omics Analysis \& Pipeline Development}: Proficient in building and deploying automated analysis pipelines for diverse genomic data modalities, including Whole Genome Sequencing (WGS), Bulk RNA-seq, single-cell RNA-seq (scRNA-seq), and Spatial Transcriptomics.
  \item \textbf{Phylogenetics \& Evolutionary Analysis}: Experience in Bayesian phylogenetic inference (\texttt{BEAST}), maximum likelihood methods (e.g., \texttt{IQ-TREE}), and population genetics simulation (\texttt{SLiM}).
  \item \textbf{Sequence Analysis}: Viral lineage assignment (\texttt{Pangolin}), sequence alignment (\texttt{BWA}, \texttt{Bowtie2}), variant calling \& manipulation (\texttt{VCFtools}, \texttt{BCFtools}), and general sequence analysis (\texttt{Samtools}, \texttt{BLAST}).
  \item \textbf{Genomic Data Formats}: Deep familiarity with standard formats including VCF, BAM/SAM, FASTQ, FASTA, and GFF/GTF.
\end{cvitems}

\vspace{4mm}
\textbf{Developer Tools \& Workflow Management}
\vspace{4.5mm}
\begin{cvitems}
  \item \textbf{Version Control}: \texttt{Git} / \texttt{GitHub} for collaborative and reproducible research.
  \item \textbf{Environment Management}: \texttt{Conda} for managing software dependencies and creating reproducible environments.
  \item \textbf{High-Performance Computing (HPC)}: Experience with job scheduling systems (e.g., \texttt{Slurm}, \texttt{SGE}) and working in a cluster environment.
  \item \textbf{Workflow Languages}: Familiarity with the principles of workflow management using \texttt{Snakemake} / \texttt{Nextflow}.
\end{cvitems}

\vspace{4mm}
\textbf{Deep Learning}
\vspace{4.5mm}
\begin{cvitems}
  \item \textbf{Machine Learning Foundations}: Strong foundation in Regression, Classification, Clustering, DNNs, and CNNs.
  \item \textbf{Advanced Sequence Models}: Understanding of advanced architectures including RNNs, LSTMs, and Transformers.
  \item \textbf{Technical Implementation}: Hands-on project experience applying \texttt{Scikit-learn}, \texttt{PyTorch}, and \texttt{TensorFlow} to bioinformatics tasks.
\end{cvitems}
\end{cvparagraph}

%-------------------------------------------------------------------------------
%  PUBLICATIONS  (newest first — same order as limengting1001.com)
%-------------------------------------------------------------------------------
\cvsection{Publications}

\begin{cvparagraph}
\begin{enumerate}[leftmargin=1.1cm, labelsep=0.5cm, align=left, itemsep=1.6mm, topsep=0pt, label=\arabic*.]
  \item Gu, H., Li, J., Sun, W., \textbf{Li, M.}, Leung, K., Wu, J. T., Yuan, H.-Y., Wang, M. H., Yang, B., McKay, M. R., Ning, N., \& Poon, L. L. M. (2026). Optimizing global genomic surveillance for early detection of emerging SARS-CoV-2 variants. \textit{Nature Communications}, 17, 4322. \url{https://doi.org/10.1038/s41467-026-70664-0}

  \item Gu, H., \textbf{Li, M.}, Wang, R., Lee, L. K., Tam, W. Y., Ho, A. Y. M., Yau, M. C. Y., Yip, K. T., Wong, R. C. W., Wong, B. K. C., Luk, K. S., Lam, J. Y. W., Que, T. L., Chow, V. C. Y., Chau, S. K. Y., Duchene, S., Siu, G. K. H., \& Poon, L. L. M. (2026). Genomic epidemiology of hospital-associated SARS-CoV-2 clusters in Hong Kong during a period of relaxed visitation (May--August 2022). \textit{Influenza and Other Respiratory Viruses}, 20(3), e70249. \url{https://doi.org/10.1111/irv.70249}

  \item Cao, Y., Chin, A. W. H., Gu, H., \textbf{Li, M.}, \& Poon, L. L. M. (2025). Long non-coding RNA BCAR4 is required for efficient influenza A virus replication. \textit{bioRxiv} [Preprint]. \url{https://doi.org/10.1101/2025.02.06.636885}

  \item Ma, Z., \textbf{Li, M.}, Li, F., Wu, K., Wu, X., Luo, T., Gao, N., Luo, H., Sui, Z., Yu, Z., Jiang, H., Shang, X., Chen, C., Yue, J., Meng, F., Duan, X., \& Xu, B. (2025). Multi-omics sequencing of gastroesophageal junction adenocarcinoma reveals prognosis-relevant key factors and a novel immunogenomic classification. \textit{Gastric Cancer}, 28(3), 344--357. \url{https://doi.org/10.1007/s10120-025-01585-y}

  \item Cao, Y., Chin, A. W. H., Gu, H., \textbf{Li, M.}, Gu, Y., Lau, S. P. N., Hui, K. P. Y., Chan, M. C. W., \& Poon, L. L. M. (2025). An interferon-stimulated long non-coding RNA USP30-AS1 as an immune modulator in influenza A virus infection. \textit{PLOS Pathogens}, 21(1), e1012854. \url{https://doi.org/10.1371/journal.ppat.1012854}

  \item Zhang, J., \& \textbf{Li, M.} (2020). Bioinformatics analysis and design based on RStudio and databases. \textit{Modern Information Technology}, 4(4), 76--79. \url{https://qikan.cqvip.com/Qikan/Article/Detail?id=7101730810}
\end{enumerate}
\end{cvparagraph}

\label{acvlastpage}
\end{document}
